\scanpage{2}
% slanted
\margin{$\mathrm{Ho}\, C$ pour $C$ une cat. à cofibrations. ---}
\margin{?}

Soit $(C,\mathrm{W}_{\mathrm e},\mathrm{cof})$ une catégorie satisfaisant C1 à C3 de Baues, (pas néc. C4) (p. ex on pourrait avoir $\mathrm{cof}=\unsure{\mathrm{cof}|\mathrm{wcof}}_{\mathrm{we}}$). \inserted{On suppose qu'il y a objet initial $\varnothing$} Alors il est connu que
\[
 W_{\mathrm{Cof}}^{-1}\,\mathrm{Cof}\longrightarrow W^{-1}C
\]
est une équivalence de catégories.

Pour le voir, on va définir un foncteur en sens inverse (qui sera quasi-inverse)
\[
 \Phi\colon W^{-1}C\longrightarrow W_{\mathrm{Cof}}^{-1}\,\mathrm{Cof},\quad
 \text{i.e.}\quad C\longrightarrow W_{\mathrm{Cof}}^{-1}.\mathrm{Cof}
 \quad\text{qui rend les }u\in W\text{ inv.}
\]

Pour tt $x\in C$, utilisant C3 pour $\varnothing\longrightarrow x$, on trouve
\[
 \widetilde x\xrightarrow[\sim]{\alpha_x}x
 \qquad\text{avec }\widetilde x\in\mathrm{Cof},\ \alpha_x\in W .
\]

Soit alors $x\xrightarrow{f}y$ dans $C$,
\diagram{\begin{tikzcd}[ampersand replacement=\&]
\widetilde x \arrow[d,"\sim"',"{\alpha_x}"'{pos=.65}] \arrow[dr,"{f\alpha_x}"] \& \widetilde y \arrow[d,"{\alpha_y}"]\\
x \arrow[r,"f"'] \& y
\end{tikzcd}}
considérons
\diagram{\begin{tikzcd}[ampersand replacement=\&]
\varnothing \arrow[d] \arrow[r] \& \widetilde y \arrow[d,"\sim"' ,"{\alpha_y}"]\\
\widetilde x \arrow[r,"{f\alpha_x}"] \& y
\end{tikzcd}}
appliquant le \lq\lq weak lifting lemma\rq\rq, on trouve
\diagram{\begin{tikzcd}[ampersand replacement=\&,column sep=small]
\& \widetilde y \arrow[r,tail,"i","\sim"'] \arrow[d,"\sim"'," {\alpha_y}"] \& \widetilde y' \arrow[dl,"\sim" ,"{\alpha'_{\deleted{\unsure{y|g}}}}"' ]\\
\widetilde x \arrow[urr,bend left=12,"{\widetilde\varphi}"] \arrow[r,"{f\alpha_x}","f"'] \& y \&
\end{tikzcd}}
un relèvement $\widetilde\varphi$ \deleted{\unsure{vers, pas de}} de $f\alpha_x$ non pas \deleted{$f\alpha_x$} vers $\widetilde y$, mais vers un $\widetilde y'$.

Considérons alors
\[
 [i]^{-1}\widetilde\varphi\colon\widetilde x\longrightarrow\widetilde y
 \qquad\text{dans }\mathrm{Hom}_{W_{\mathrm{Cof}}^{-1}}(\widetilde x,\widetilde y)
\]
Je dis qu'il ne dépend pas du choix du $(i,\alpha'_{\unsure{y|g}},\varphi)$. Pour deux tels choix \deleted{\unsure{on}}
\diagram{\begin{tikzcd}[ampersand replacement=\&]
\& \widetilde y'_1\\
\widetilde y \arrow[ur,tail,"{i_1}"] \arrow[dr,tail,"{i_2}"'] \&\\
\&\widetilde y'_2
\end{tikzcd}}
parallèles : introduire $\widetilde y'_3=\widetilde y'_1\vee_{\widetilde y}\widetilde y'_2$, OPS que $\widetilde y'_1=\widetilde y'_2$, $i_1=i_2$, $\alpha'_1=\alpha'_2$, donc on a deux relèvements $\varphi_1,\varphi_2$ de $f\alpha_x$ vers la $\mathord{=}$ $\widetilde y'$, \deleted{OPS \unsure{$=$|et} que $\widetilde y$} et on est ramené (écrivant $y'$ au lieu de $\widetilde y'$, $\deleted{\unsure{g}}\inserted{g=f\alpha_x}$, $x'$ au lieu de $\widetilde x$, mais $x$ cof.)

\diagram{\begin{tikzcd}[ampersand replacement=\&]
x' \arrow[r,shift left,"\varphi"] \arrow[r,shift right,"\psi"'] \& y' \arrow[r,"\alpha","\sim"'] \& y
\end{tikzcd}}
\deleted{\diagram{\begin{tikzcd}[ampersand replacement=\&]
\& y'\ {\rm cof}\arrow[d,"\alpha","\sim"']\\
x'\ {\rm cof}\arrow[ur,shift left,"\varphi"]\arrow[ur,shift right,"\psi"']\arrow[r,"g"']\&y
\end{tikzcd}}}
prouver que si $x',y'\in\mathrm{Cof}$, et $\alpha\varphi=\alpha\psi\deleted{\unsure{{}=g|\alpha g}}$, alors $[\varphi]=[\psi]$ dans $W_{\mathrm{Cof}}^{-1}\,\mathrm{Cof}$.

\diagram{\begin{tikzcd}[ampersand replacement=\&]
x'\arrow[r]\& y\\
x'\amalg x'\arrow[u,"{\nabla_{x'}}"]\arrow[r,"{(\varphi,\psi)}"']\& y'\arrow[u,"{\in W}"]
\end{tikzcd}}
\diagram{\begin{tikzcd}[ampersand replacement=\&]
\varnothing\arrow[r]\&\unsure{x|{\cdot}}\arrow[d]\\
\&y
\end{tikzcd}\qquad
\begin{tikzcd}[ampersand replacement=\&]x'_{\rm cof}\arrow[r,shift left]\arrow[r,shift right]\&y'_{\rm cof}\arrow[r,line width=1.3pt]\&y\end{tikzcd}}
Si ce lemme est valable, on trouve un foncteur $\deleted{\unsure{W}}C\longrightarrow W_{\mathrm{Cof}}^{-1}\,\mathrm{Cof}$, car la compatibilité de $u\mapsto\Phi(u)$ aux identités et aux compositions est immédiate.
% archive sheet number: 1
